\xchapter{The Thirty-Nine Articles of Religion, A.D. 1571}

\emph{Together with the Revision of the Same, as set forth by the}

PROTESTANT EPISCOPAL CHURCH IN THE UNITED STATES OF AMERICA,

A.D. 1801.

[1. The \emph{Latin} text of the Elizabethan Articles, adopted in 1562, is a reprint of the \emph{editio princeps} of Reginald Wolfe, royal printer, London, 1563, issued by express authority of the Queen, and reproduced by Charles Hardwick, in his \emph{History of the Articles of Religion}, new edition, Cambridge, 1859, pp. 277 sqq. (Hardwick gives also, in four parallel columns, the English edition of 1571, and the Forty-two Articles of 1553, Latin and English, with the textual variations of the Parker MS. of 1571, and other printed editions.)

2. The English text is reprinted, with the old spelling, from the authorized London edition of John Cawood, 1571, as found in Hardwick, l. c.

The question of the comparative authority of the Latin and English texts is answered by Burnet, Waterland, and Hardwick, to the effect that both are equally authentic, but that in doubtful cases the Latin must determine the sense. The Articles were passed, recorded, and ratified to the year 1562 (1563), in \emph{Latin} only; but these Latin Articles were revised and translated by the Convocation of 1571, and both the Latin and English texts, adjusted as nearly as possible, were published in the same year by the royal authority. Subscription was hereafter required to the English Articles, called the Articles of 1562, by the famous Act of the XIII. of Elizabeth. See Hardwick, l. c. p. 159.

3. The \emph{American} Revision of the Articles, as adopted by the General Convention of the Protestant Episcopal Church in the United States, held in Trenton, New Jersey, Sept. 12, 1801, is taken from the standard American edition of \emph{The Book of Common Prayer} (published by the Harpers, New York, 1844, and by the New York Bible and Common Prayer-Book Society, 1873, pp. 512 sqq.). It has been compared with the Journal of the Convention, edited by Dr. W. S\textsc{tevens} P\textsc{erry}, in \emph{Journals of the General Convention of the Protestant Episcopal Church in the United States}, 1785–1835 (Claremont, N. H., 1874), Vol. I. pp.279 sqq.

4. To facilitate the comparison, the words in which the English and American editions differ are printed in italics. The chief differences are the omission of the Athanasian Creed, in Art. VIII.; the omission of Art. XXI., on the Authority of General Councils; and the entire reconstruction of Art. XXXVII., on the Power of the Civil Magistrate.]

The English editions of the Articles are usually preceded by the following Royal Declaration, which is the work of Archbishop Laud (1628):

'Being by God's Ordinance, according to Our just Title, \emph{Defender of the Faith, and Supreme Governour of the Church, within these Our Dominions}, We hold it most agreeable to this Our Kingly Office, and Our own religious Zeal, to conserve and maintain the Church committed to Our Charge, in the Unity of true Religion, and in the Bond of Peace; and not to suffer unnecessary Disputations, Altercations, or Questions to be raised, which may nourish Faction both in the Church and Commonwealth. We have therefore, upon mature Deliberation, and with the Advice of so many of Our Bishops as might conveniently be called together, thought fit to make this Declaration following:

'That the Articles of the Church of \emph{England} (which have been allowed and authorized heretofore, and which Our Clergy generally have subscribed unto) do contain the true Doctrine of the \emph{Church of England} agreeable to God's Word, which We do therefore ratify and confirm, requiring all Our loving Subjects to continue in the uniform Proffession thereof, and prohibiting the least difference from the said Articles; which to that End We command to be new printed, and this Our Declaration to be published therewith.

'That We are Supreme Governour of the Church of \emph{England: } And that if any Difference arise about the external Policy, concerning the \emph{Injunctions, Canons}, and other \emph{Constitutions} whatsoever thereto belonging, the Clergy in their Convocation is to order and settle them, having first obtained leave under Our Broad Seal so to do: and We approving their said Ordinances and Constitutions, providing that none be made contrary to the Laws and Customs of the Land.

'That out of Our Princely Care that the Churchmen may do the Work which is proper unto them, the Bishops and Clergy, from time to time in Convocation, upon their humble Desire, shall have Licence under Our Broad Seal to deliberate of, and to do all such Things, as, being made plain by them, and assented unto by iis, shall concern the settled Continuance of the Doctrine and Discipline of the Church of \emph{England} now established; from which We will not endure any varying or departing in the least Degree.

'That for the present, though some differences have been ill raised, yet We take comfort in this, that all Clergymen within Our Realm have always most willingly subscribed to the Articles established; which is an argument to Us, that they all agree in the true, usual, literal meaning of the said Articles; and that even in those curious points, in which the present differences lie, men of all sorts take the Articles of the Church of \emph{England} to be for them; which is an argument again, that none of them intend any desertion of the Articles established.

'That, therefore, in these both curious and unhappy differences, which have for so many hundred years, in different times and places, exercised the Church of Christ, We will, that all further curious search be laid aside, and these disputes shut up in God's promises, as they be generally set forth to us in the holy Scriptures, and the general meaning of the Articles of the Church of England according to them. And that no man hereafter shall either print or preach to draw the Article aside any way, but shall submit to it in the plain and full meaning thereof: and shall not put his own sense or comment to be the meaning of the Article, but shall take it in the literal and grammatical sense.

'That if any publick Reader in either of Our Universities, or any Head or Master of a College, or any other person respectively in either of them, shall affix any new sense to any Article, or shall publickly read, determine, or hold any publick Disputation, or suffer any such to be held either way, in either the Universities or Colleges respectively; or if any Divine in the Universities shall preach or print any thing either way, other than is already established in Convocation with our Royal Assent, he, or they the Offenders, shall be liable to our displeasure, and the Church's censure in Our Commission Ecclesiastical, as well as any other: And we will see there shall be due Execution upon them.'

E\textsc{ditio} L\textsc{atina} P\textsc{rinceps},

E\textsc{nglish} E\textsc{dition},

A\textsc{merican} R\textsc{evision},

1563 [1562].

1571.

1801.

\emph{Articuli, de quibus in Synodo Londinensi anno Domini, iuxta ecclesiæ Anglicanæ computationem, M.D.LXII. ad tollendam opinionum, dissensionem, et firmandum in uera Religione consensum, inter Archiepiscopos Episcoposque utriusque Prouinciæ, nec non etiam uniuersum Clerum convenit.}

Articles whereupon it was agreed by the Archbishoppes and Bishoppes of both prouinces and the whole cleargie, in the Conuocation holden at London in the yere of our Lorde God, 1562. according to the computation of the Churche of Englande, for the auoiding of the diuersities of opinions, and for the stablishyng of consent touching true Religion.

Articles of Religion; as established by the Bishops, the Clergy, and Laity of the Protestant Episcopal Church in the United States of America, in Convention, on the twelfth day of September, in the year of our Lord 1801.

\xsection{Article I. Of Faith in the Holy Trinity}

\emph{Vnus est viuus et uerus Deus æternus, incorporeus, impartibilis, impassibilis, immensæ potentiæ, sapientiæ ac bonitatis: creator et conseruator omnium tum uisibilium tum inuisibilium. Et in Vnitate huius diuinæ naturæ tres sunt Personæ, }

There is but one lyuyng and true God, euerlastyng, without body, partes, or passions, of infinite power, wysdome, and goodnesse, the maker and preseruer of al things both visible and inuisible. And in vnitie of this Godhead

There is but one living and true God, everlasting, without body, parts, or passions; of infinite power, wisdom, and goodness; the Maker, and Preserver of all things both visible and invisible. And in unity of this God-head there be

\emph{eiusdem essentiæ, potentiæ, ac æternitatis, Pater, Filius, et Spiritus sanctus.}

there be three persons, of one substaunce, power, and eternitie, the father, the sonne, and the holy ghost.

three Persons, of one substance, power, and eternity: the Father, the Son, and the Holy Ghost.

\xsection{Article II. Of the Word or Son of God, which was made very Man}

\emph{Filius, qui est uerbum Patris ab æterno à Patre genitus uerus et æternus Deus, ac Patri consubstantialis, in utero Beatæ uirginis ex illlus substantia naturam humanam assumpsit: ita ut duæ naturæ, diuina et humana integrè atque perfectè in unitate personæ, fuerint inseparabiliter coniunctæ: ex quibus est vnus CHRISTVS, verus Deus et verus Homo: qui uerè passus est, crucifixus, mortuus, et sepultus, ut Patrem nobis reconciliaret, essetque} [\emph{hostia}] \emph{non tantùm pro culpa originis, uerum etiam pro omnibus actualibus hominum peccatis.}

The Sonne, which is the worde of the Father, begotten from euerlastyng of the Father, the very and eternall GOD, of one substaunce with the Father, toke man's nature in the wombe of the blessed Virgin, of her substaunce: so that two whole and perfect natures, that is to say, the Godhead and manhood, were ioyned together in one person, neuer to be diuided, whereof is one Christe, very GOD and very man, who truely suffered, was crucified, dead, and buried, to reconcile his father to vs, and to be a sacrifice, not only for originall gylt, but also for \emph{all}\footnote{The omission of '\emph{all}' dates from the year 1630, and the revised text of the Westminster Assembly of Divines, 1647. It appears in the edition of 1628, and is restored in modern English editions. See Hardwick, p. 279.} actuall sinnes of men.

The Son, which is the Word of the Father, begotten from everlasting of the Father, the very and eternal God, \emph{and} of one substance with the Father, took Man's nature in the womb of the blessed Virgin, of her substance: so that two whole and perfect Natures, that is to say, the Godhead and Manhood, were joined together in one Person, never to be divided, whereof is one Christ, very God, and very Man; who truly suffered, was crucified, dead, and buried, to reconcile his Father to us, and to be a sacrifice, not only for original guilt, but also for actual sins of men.

\xsection{Article III. Of the going down of Christ into Hell}

\emph{Qvemmadmodum Christus pro nobis mortuus est et sepultus, ita est etiam credendus ad Inferos descendisse.}

As Christe dyed for vs, and was buryed: so also it is to be beleued that he went downe into hell.

As Christ died for us, and was buried, so also \emph{is it} to be believed that he went down into Hell.

\xsection{Article IV. Of the Resurrection of Christ}

Christus vere a mortuis resurrexit, suumque corpus cum carne, ossibus, omnibusque ad integritatem humanæ naturæ pertinentibus, recepit, cum quibus in cœlum ascendit, ibique residet, quoad extremo die ad iudicandos [omnes] homines reuersurus sit.

Christe dyd truly \emph{aryse} agayne from death, and toke agayne his body, with flesh, bones, and all thinges apperteyning to the perfection of man's nature, wherewith he ascended into heauen, and there sitteth, vntyll he returne to iudge all men at the last day.

Christ did truly \emph{rise} again from death, and took again his body, with flesh, bones, and all things appertaining to the perfection of Man's nature; wherewith he ascended into Heaven, and there sitteth, until he return to judge all Men at the last day.

\xsection{Article V. Of the Holy Ghost}

Spiritus sanctus, à patre et filio procedens, eiusdem est cum patre et filio essentiæ, maiestatis, et gloriæ, uerus, ac æternus Deus.

The holy ghost, proceedyng from the Father and the Sonne, is of one substaunce, maiestie, and glorie, with the Father, and the Sonne, very and eternall God.

The Holy Ghost, proceeding from the Father and the Son, is of one substance, majesty, and glory, with the Father and the Son, very and eternal God.

\xsection{Article VI. Of the Sufficiency of the Holy Scriptures for Saltation}

Scriptura sacra continet omnia quæ sunt ad salutem necessaria, ita ut quicquid in ea nec legitur, neque inde probari potest, non sit à quoquam exigendum, ut tanquam Articulus fidei credatur, aut ad necessitatem salutis requiri putetur.

Holy Scripture conteyneth all thinges necessarie to saluation: so that whatsoeuer is not read therein, nor may be proued therby, is not to be required of anye man, that it shoulde be beleued as an article of the fayth, or be thought requisite [as] necessarie to saluation.

Holy Scripture containeth all things necessary to salvation: so that whatsoever is not read therein, nor may be proved thereby, is not to be required of any man, that it should be believed as an article of the Faith, \emph{or} be thought requisite or necessary to salvation.

Sacræ Scripturæ nomine eos Canonicos libros Veteris et Novi testamenti intelligimus,

In the name of holy Scripture, we do vnderstande those Canonicall

In the name of \emph{the} Holy Scripture we do understand those canonical Books of

de quorum autoritate in Ecclesia, nunquam dubitatum est.

bookes of the olde and newe Testament, of whose aucthoritie was neuer any doubt in the Churche.

the Old and New Testament, of whose authority was never any doubt in the Church.

Catalogus librorum sacræ Canonicæ scripturæ Veteris Testamenti.

Of the names and number of the Canonicall Bookes.

Of the Names and Number of the Canonical Books.

\emph{Genesis}.

Genesis.

Genesis,

\emph{Exodus.}

Exodus.

Exodus,

\emph{Leviticus.}

Leuiticus.

Leviticus,

\emph{Numeri.}

Numerie.

Numbers,

\emph{Deuteronom.}

Deuteronomium.

Deuteronomy,

\emph{Iosue.}

Iosue.

Joshua,

\emph{Iudicum.}

Iudges.

Judges,

\emph{Ruth.}

Ruth.

Ruth,

\emph{2 Regum.}

The .1. boke of Samuel

The First Book of Samuel,

\emph{Paralipom. 2.}

The .2. boke of Samuel.

The Second Book of Samuel,

\emph{2 Samuelis.}

The .1. booke of Kinges.

The First Book of Kings,

The .2. booke of Kinges.

The Second Book of Kings,

The .1. booke of Chroni.

The First Book of Chronicles,

The .2. booke of Chroni.

The Second Book of Chronicles,

\emph{Esdræ.}

The .1. booke of Esdras.

The First Book of Esdras,

The .3. booke of Esdras.

The Second Book of Esdras,

\emph{Hester.}

The booke of Hester.

The Book of Esther,

\emph{Iob.}

The booke of Iob.

The Book of Job,

\emph{Psalmi.}

The Psalmes.

The Psalms,

\emph{Prouerbia.}

The Prouerbes.

The Proverbs,

\emph{Ecclesiastes.}

Ecclesia, or preacher.

Ecclesiastes or Preacher,

\emph{Cantica.}

Cantica, or songes of Sa.

Cantica, or Songs of Solomon,

\emph{Prophetæ maiores.}

4. Prophetes the greater.

Four Prophets the greater,

\emph{Prophetæ minores.}

12. Prophetes the lesse.

Twelve Prophets the less.

Alios autem Libros (ut ait Hieronymus) legit quidem Ecclesia ad exempla uitæ et formandos mores, illos tamen

And the other bookes, (as Hierome sayth) the Churche doth reade for example of lyfe and instruction

And the other Books (as Hierome saith) the Church doth read for example of life and instruction of manners:

ad dogmata confirmanda non adhibet: ut sunt

of manners: but yet doth it not applie them to establishe any doctrene. Such are these followyng.

but yet doth it not apply to them to establish any doctrine: such are these following:

\emph{Tertius et quartus Esdræ.}

The third boke of Esdras.

The Third Book of Esdras,

\emph{Sapientia.}

The fourth boke of Esdras.

The Fourth Book of Esdras,

\emph{Iesus filius Syrach.}

The booke of Tobias.

The Book of Tobias,

\emph{Tobias.} .\emph{Iudith.}

The booke of Iudith.

The Book of Judith,

\emph{Libri Machabæorum. 2.}

The reste of the booke of Hester.

The rest of the Book of Esther,

The booke of Wisdome.

The Book of Wisdom,

Iesus the sonne of Sirach.

Jesus the Son of Sirach,

Baruch the prophet.

Baruch the Prophet,

Song of the .3. Children.

Song of the Three Children,

The storie of Susanna.

The Story of Susanna,

Of Bel and the Dragon.

Of Bel and the Dragon,

The prayer of Manasses.

The prayer of Manasses,

The .1. boke of Machab.

The First Book of Macabees,

The .2. Booke of Macha.

The Second Book of Macabees,

Noui Testamenti Libros omnes (ut uulgo recepti sunt) recipimus et habemus pro Canonicis.

All the bookes of the newe Testament, as they are commonly receaued, we do receaue and accompt them \emph{for} Canonicall.

All the Books of the New Testament, as they are commonly received, we do receive, and account them Canonical.

\xsection{Article VII. Of the Old Testament}

Testamentum vetus Nouo contrarium non est, quandoquidem tam in veteri quam nova, per Christum, qui vnicas est mediator Dei et hominum, Deus et Homo, æterna vita humano generi est proposita. Quare malè sentiunt,

The olde Testament is not contrary to the newe, for both in the olde and newe Testament euerlastyng lyfe is offered to mankynde by Christe, who is the onlye mediatour betweene God and man.

The Old Testament is not contrary to the New: for both in the Old and New Testament everlasting life is offered to Mankind by Christ, who is the only Mediator between God and Man, being both God and

qui veteres tantum in promissiones temporarias sperasse confingunt. Quanquam Lex à Deo data per Mosen, quoad Ceremonias et ritus, Christianos non astringat, neque ciuilia eius præcepta in aliqua Republica necessariò recipi debeant: nihilominus tamen ab obedientia mandatorum, quæ Moralia vocantur, nullus quantumius Christianus, est solutus.

Wherefore they are not to be hearde whiche faigne that the olde fathers dyd looke only for transitorie promises. Although the lawe geuen from God by Moses, as touchyng ceremonies and rites, do not bynde Christian men, nor the ciuile preceptes thereof, ought of necessitie to be receaued in any common wealth: yet notwithstandyng, no Christian man whatsoeuer, is free from the obedience of the commaundementes, which are called morall.

Man. Wherefore they are not to be heard, which feign that the old Fathers did look only for transitory promises. Although the Law given from God by Moses, as touching Ceremonies and Rites, do not bind Christian men, nor the Civil precepts thereof ought of necessity to be received in any commonwealth; yet notwithstanding, no Christian man whatsoever is free from the obedience of the Commandments which are called Moral.

\xsection{Article VIII. Of the Creeds}

Symbola tria, Nicænum, Athanasij, et quod vulgo Apostolicum appellatur, omnino recipienda sunt et credenda. Nam firmissimis Scripturarum testimonijs probari possunt.

The \emph{three Credes}, Nicene Crede, \emph{Athanasian Crede}, and that whiche is commonlye called the Apostles' Crede, ought \emph{throughlye} to be receaued and beleued: for they may be proued by moste certayne warrauntes of holye scripture.

The Nicene Creed, and that which is commonly called the Apostles' Creed, ought \emph{thoroughly} to be received and believed: for they may be proved by most certain warrants of Holy Scripture.

\xsection{Article IX. Of Original or Birth-Sin}

Peccatum originis non est (et fabulantur Pelagiani) in imitatione Adami situm, sed est vitium et deprauatio naturæ cuiuslibet hominis ex Adamo naturaliter propagati, qua fit, vt ab originali iustitia quàm longissime

Originall sinne standeth not in the following of Adam (as the Pelagians do vaynely talke) but it is the fault and corruption of the nature of euery man, that naturally is engendered of the ofspring

Original sin standeth not in the following of Adam (as the Pelagians do vainly talk); but it is the fault and corruption of the Nature of every man, that naturally is engendered of the offspring

distet, ad malum sua natura propendeat, et caro semper aduersus spiritum concupiscat. Vnde in vno quoque nascentium, iram Dei atque damnationem meretur. Manet etiam in renatis hæc naturæ deprauatio; qua fit, ut affectus carnis, græce \textgreek{σαρκός}, (quod alij sapientiam, alij sensum, alij affectum, alij studium  [carnis] interpretantur) legi Dei non subjiciatu. Et quanquàm renatis et credentibus nulla propter Christum est condemnatio, peccati tamen in sese rationem habere Concupiscentiam fatetur Apostolus.

of Adam, whereby man is very farre gone from originall ryghteousness, and is of his owne nature enclined to euyll, so that the fleshe lusteth alwayes contrary to the spirite, and therefore in euery person borne into this worlde, it deserueth Gods wrath and damnation. And this infection of nature doth remayne, yea in them that are regenerated, whereby the luste of the fleshe, called in Greke \textgreek{φρόνημα σαρκὸς}, which some do expounde the wisdome, some sensualitie, some the affection, some the desyre of the fleshe, is not subiect to the lawe of God. And although there is no condemnation for them that beleue and are baptized: yet the Apostle doth confesse that concupiscence and luste hath of it selfe the nature of synne.

of Adam; whereby man is very far gone from original righteousness, and is of his own nature inclined to evil, so that the flesh lusteth always contrary to the spirit; and therefore in every person born into this world, it deserveth God's wrath and damnation. And this infection of nature doth remain, yea in them that are regenerated; whereby the lust of the flesh, called in Greek \textgreek{φρόνημα σαρκὸς} (which some do expound the wisdom, some sensuality, some the affection, some the desire, of the flesh), is not subject to the Law of God. And although there is no condemnation for them that believe and are baptized; yet the Apostle doth confess, that concupiscence and lust hath of itself the nature of sin.

\xsection{Article X. Of Free-Will}

Ea est hominis post lapgum Adæ conditio, ut sese naturalibus suis viribus et bonis operibus ad fidem et invocationem Dei conuertere ac præparare non possit: Quare absque gratia Dei, quæ per Christum est, nos præueniente, ut uelimus, et cooperante dum volumus, ad

The condition of man after the fall of Adam is suche, that he can not turne and prepare hym selfe by his owne naturall strength and good workes, to fayth and calling vpon God: Wherefore we haue no power to do good workes pleasaunt and acceptable

The condition of Man after the fall of Adam is such, that he can not turn and prepare himself, by his own natural strength and good works, to faith, and calling upon God. Wherefore we have no power to do good works pleasant and acceptable to

pietatis opera facienda, quæ Deo grata sint et accepta, nihil valemus.

to God, without the grace of God by Christe preuentyng us, that we may haue a good wyll, and workyng with vs, when we haue that good wyll.

God, without the grace of God by Christ preventing us, that we may have a good will, and working with us, when we have that good will.

\xsection{Article XI. Of the Justification of Man}

Tantùm propter meritum Domini ac Seruatoris nostri Iesu Christi, per fidem, non propter opera et merita nostra, iusti coram Deo reputamur: Quare sola fide nos iustificari, doctrina est saluberrima, ac consolationis plenissima: ut in Homilia de Iustificatione hominis fusiùs explicatur.

We are accompted righteous before God, only for the merite of our Lord and sauiour Jesus Christe, by faith, and not for our owne workes or deseruynges. Wherefore, that we are iustified by fayth onely, is a most wholesome doctrine, and very full of comfort, as more largely is expressed in the Homilie of iustiflcation.

We are accounted righteous before God, only for the merit of our Lord and Saviour Jesus Christ by Faith, and not for our own works or deservings. Wherefore, that we are justified by Faith only, is a most wholesome Doctrine, and very full of comfort, as more largely is expressed in the Homily of Justification.

\xsection{Article XII. Of Good Works}

Bona opera quæ sunt fructus fidei et iustificatos sequuntur, quanquam peccata nostra expiari et diuini iudicij seueritatem ferre non possunt, Deo tamen grata sunt et accepta in Christo, atque ex uera et uiua fide necessario profluunt, ut plane ex illis, æque fides uiua cognosci possit, atque arbor ex fructu iudicari.

Albeit that good workes which are the fruites of fayth, and folowe after iustification, can not put away our sinnes, and endure the seueritie of Gods iudgement: yet are they pleasing and acceptable to God in Christe, and do spring out necessarily of a true and liuely fayth, in so muche that by them, a lyuely fayth may be as euidently knowen, as a tree discerned by the fruit.

Of Good Works. Albeit that Good Works, which are the fruits of Faith, and follow after Justification, can not put away our sins, and endure the severity of God's judgment; yet are they pleasing and acceptable to God in Christ, and do spring out necessarily of a true and lively Faith; insomuch that by them a lively Faith may be as evidently known as a tree discerned by the fruit.

\xsection{Article XIII. Of Works before Justification}

\emph{Opera quæ fiunt ante gratiam Christi, et spiritus eius afflatum, cum ex fide Iesu Christi non prodeant, minimè Deo grata sunt: neque gratiam }(\emph{ut multi uocant}) \emph{de congruo merentur: Imo cum non sint facta ut Deus illa fieri uoluit et præcepit, peccati rationem habere non dubitamus.}

Workes done before the grace of Christe, and the inspiration of his spirite, are not pleasaunt to God forasmuche as they spring not of fayth in Jesu Christ, neither do they make men meete to receaue grace, or (as the schole aucthours saye) deserue grace of congruitie: yea rather for that they are not done as GOD hath wylled and commaunded them to be done, we doubt not but they haue the nature of synne.

Works done before the grace of Christ, and the Inspiration of his Spirit, are not pleasant to God, forasmuch as they spring not of faith in Jesus Christ; neither do they make men meet to receive grace, or (as the School-authors say) deserve grace of congruity: yea rather, for that they are not done as God hath willed and commanded them to be done, we doubt not but they have the nature of sin.

\xsection{Article XIV. Of Works of Supererogation}

Opera quæ supererogationis appellant, non possunt sine arrogantia et impietate prædicari. Nam illis declarant homines non tantum se Deo reddere quæ tenentur sed plus in eius gratiam facere quam deberent: cum apertè Christus dicat: Cum feceritis omnia quæcunque præcepta sunt uobis, dicite: Serui inutiles sumus.

Voluntarie workes besydes, ouer and aboue Gods commaundementes, which they call workes of supererogation, can not be taught without arrogancie and impietie. For by them men do declare that they do not onely render vnto God as muche as they are bounde to do, but that they do more for his sake than of bounden duetie is required: Whereas Christe sayth playnly, When ye have done al that are commaunded to you, say, We \emph{be} vnprofltable seruantes.

Voluntary Works besides, over and above, God's Commandments, which they call Works of Supererogation, can not be taught without arrogancy and impiety: for by them men ao declare, that they do not only render unto God as much as they are bound to do, but that they do more for his sake, than of bounden duty is required: whereas Christ saith plainly, When ye have done all that are commanded to you, say, We \emph{are} unprofitable servants.

\xsection{Article XV. Of Christ alone without Sin}

\emph{Christus in nostræ naturæ ueritate per omnia similis factus est nobis, excepto peccato, à quo prorsus erat immunis, tum in carne tum in spiritu. Venit, ut Agnus absque macula esset, qui mundi peccata per immolationem sui semel factam, tolleret: et peccatum} (\emph{ut inquit Ioannes}) \emph{in eo non erat. Sed nos reliqui, etiam baptizati, et in Christo regenerati, in multis tamen offendimus omnes: Et si dixerimus quia peccatum non habemus, nos ipsos seducimus, et veritas in nobis non est.}

Christe in the trueth of oure nature, was made lyke vnto vs in al thinges (sinne only except) from which he was clearley voyde, both in his fleshe, and in his spirite. He came to be the lambe without spot, who by \emph{the} sacrifice of hym self once made, shoulde take away the sinnes of the worlde: and sinne, (as S. John sayeth) was not in hym. But al we the rest, (although baptized, and borne agayne in Christ) yet offende in many thinges, and if we say we haue no sinne, we deceaue our selues, and the trueth is not in vs.

Christ in the truth of our nature was made like unto us in all things, sin only except, from which he was clearly void, both in his flesh, and in his spirit. He came to be the Lamb without spot, who, by sacrifice of himself once made, should take away the sins of the world; and sin (as Saint John saith) was not in him. But all we the rest, although baptized, and born again in Christ, yet offend in many things; and if we say we have no sin, we deceive ourselves, and the truth is not in us.

\xsection{Article XVI. Of Sin after Baptism}

Non omne peccatum Mortale post baptismum uoluntarie perpetratum, est peccatum in Spiritum sanctum et irremissible. Proinde lapsis à baptismo in peccata, locus pœnitentia non est negandus. Post acceptum spiritum sanctum, possumus à gratia data recedere atque peccare, denuóque per gratiam Dei resurgere ac respiscere. Ideóque illi damnandi sunt, qui se quamdiu

Not euery deadly sinne willingly committed after baptisme, is sinne agaynst the holy ghost, and vnpardonable. Wherefore the graunt of repentaunce is not to be denyed to such as fall into sinne after baptisme. After we have receaued the holy ghost, we may depart from grace geuen, and fall into sinne, and by the grace of God (we may) aryse agayne,

Not every deadly sin willingly committed after Baptism is sin against the Holy Ghost, and unpardonable. Wherefore the grant of repentance is not to be denied to such as fall into sin after Baptism. After we have received the Holy Ghost, we may depart from grace given, and fall into; sin, and by the grace of God we may arise again, and amend our lives. And

hic viuant, amplius non posse peccare affirmant, aut verè resipiscentibus pœnitentiæ locum denegant.

and amend our lyues. And therefore, they are to be condemned, whiche say they can no more sinne as long as they lyue here, or denie the place of forgeuenesse to such as truely repent.

therefore they are to be condemned, which say, they can no more sin as long as they live here, or deny the place of forgiveness to such as truly repent.

\xsection{Article XVII. Of Predestination and Election}

Prædestinatio ad uitam, est æternum Dei propositum, quo ante iacta mundi fundamenta, suo consilio, nobis quidem occulto, constanter decreuit, eos quos in Christo elegit ex hominum genere, à maledicto et exitio liberare, atque ut uasa in honorem efficta, per Christum ad æternam salutem adducere: Vnde qui tam præclaro Dei beneficio sunt donati, illi spiritu eius opportuno tempore operante, secundum propositum eius uocantur: uocationi per gratiam parent: iustificantur gratis: adoptantur in filios; vnigeniti Iesu Christi imagini efficiuntur conformes: in bonis operibus sanctè ambulant: et demùm ex Dei misericordia pertingunt ad sempiternam fœlicitatem.

Predestination to lyfe, is the euerlastyng purpose of God, whereby (before the foundations of the world were layd) he hath constantly decreed by his councell secrete to vs, to deliuer from curse and damnation, those whom he hath chosen in Christe out of rnankynd, and to bryng them by Christe to euerlastyng saluation, as vessels made to honour. Wherefore they which be indued with so excellent a benefite of God, be called accordyng to Gods purpose by his spirite workyng in due season: they through grace obey the callyng: they be iustified freely: they be made sonnes of God by adoption: they be made lyke the image of his onelye begotten sonne Jesus Christe: they walke religiously in good workes, and at length by Gods mercy, they attaine to euerlastyng felicitie.

Predestination to Life is the everlasting purpose of God, whereby (before the foundations of the world were laid) he hath constantly decreed by his counsel secret to us, to deliver from curse and damnation those whom he hath chosen in Christ out of mankind, and to bring them by Christ to everlasting salvation, as vessels made to honour. Wherefore, they which be endued with so excellent a benefit of God, be called according to God's purpose by his Spirit working in due season: they through Grace obey the calling: they be justified freely: they be made sons of God by adoption: they be made like the image of his only-begotten Son Jesus Christ: they walk religiously in good works, and at length, by God's mercy, they attain to everlasting felicity.

Quemadmodum Prædestinationis et Electionis nostræ in Christo pia consideratio, dulcis, suauis et ineffabilis comolationis plena est verè pijs et his qui sentiunt in se uim spiritus CHRISTI, facta carnis et membra quæ adhuc sunt super terram mortificantem, animumque ad cœlestia et superna rapientem, tum quia fidem nostram de æterna salute consequenda per Christum plurimum stabilit atque confirmat, tum quia amorem nostrum in Deum uehementer accendit; ita hominibus curiosis, carnalibus, et spiritu Christi destitutis, ob oculos perpetuò versari Prædestinationis Dei sententiam, pernitiosissimum, est præcipitium, unde illos Diabolus protrudit, uel in desperationem, uel in æquè pernitiosam impurissimæ vitæ securitatem.

As the godly consyderation of predestination, and our election in Christe, is full of sweete, pleasaunt, and vnspeakeable comfort to godly persons, and such as feele in themselues the working of the spirite of Christe, mortifying the workes of the fleshe, and their earthlye members, and drawing vp their mynde to hygh and heauenly thinges, as well because it doth greatly establyshe and confirme their fayth of eternal saluation to be enjoyed through Christe, as because it doth feruently kindle their loue towardes God. So, for curious and carnal persons, lacking the spirite of Christe, to haue continually before their eyes the sentence of Gods predestination, is a most daungerous downefall, whereby the deuyll doth thrust them either into desperation, or into rechelesnesse of most vncleane liuing, no lesse perilous then desperation.

As the godly consideration of Predestination, and our Election in Christ, is full of sweet, pleasant, and unspeakable comfort to godly persons, and such as feel in themselves the working of the Spirit of Christ, mortifying the works of the flesh, and their earthly members, and drawing up their mind to high and heavenly things, as well because it doth greatly establish and confirm their faith of eternal Salvation to be enjoyed through Christ, as because it doth fervently kindle their love towards God: So, for curious and carnal persons, lacking the Spirit of Christ, to have continually before their eyes the sentence of God's Predestination, is a most dangerous downfall, whereby the Devil doth thrust them either into desperation, or into wretchlessness of most unclean living, no less perilous than desperation.

Deinde promissiones diuinas sic amplecti oportet, ut nobis in Sacris literis generaliter propositæ sunt: ut Dei voluntas in nostris actionibus ea sequenda est,

Furthermore,\footnote{In the Forty-two Articles of 1553 there is the addition: 'Although the decrees of predestination are unknown unto us.'} we must receaue Gods promises in such wyse, as they be generally set foorth to vs in holy scripture: and in our doynges, that wyl of God is to

Furthermore, we must receive God's promises in such wise, as they be generally set forth to us in Holy Scripture; and, in our doings, that Will of God is

quam in uerbo Dei habemus disertè reuelatam.

be folowed, which we haue expreslye declared vnto vs in the worde of God.

to be followed, which we have expressly declared unto us in the Word of God.

\xsection{Article XVIII. Of obtaining eternal Salvation only by the Name of Christ}

Svnt illi anathematizandi qui dicere audent, vnumquemque in Lege aut secta quam profitetur, esse seruandum: modo iuxta illam et lumen naturæ accurate vixerit: cùm sacræ literæ tantum Iesu Christi nomen prædicent, in quo saluos fieri homines oporteat.

They also are to be had accursed, that presume to say, that euery man shal be saued by the lawe or sect which he professeth, so that he be diligent to frame his life accordyng to that lawe, and the light of nature. For holy scripture doth set out vnto vs onely the name of Jesus Christe, whereby men must be saved.

They also are to be had accursed that presume to say, That every man shall be saved by the Law or Sect which he professeth, so that he be diligent to frame his life according to that Law, and the light of Nature. For Holy Scripture doth set out unto us only the Name of Jesus Christ, whereby men must be saved.

\xsection{Article XIX. Of the Church}

Ecclesia Christi visibilis. est cœtus fidelium, in quo uerbum Dei purum prædicatur, et sacramenta, quoad ea quæ necessario exiguntur, iuxta Christi institutum recte administrantur.

The visible Church of Christe, is a congregation of faythfull men in the which the pure worde of God is preached, and the Sacramentes be duely ministred, accordyng to Christes ordinaunce in all those thynges that of necessitie are requisite to the same.

The visible Church of Christ is a congregation of faithful men, in the which the pure Word of God is preached, and the Sacraments be duly ministered according to Christ's ordinance, in all those things that of necessity are requisite to the same.

Sicut errauit ecclesia Hierosolymitana, Alexandrina et Antiochena: ita et errauit Ecclesia Romana, non solùm quoad agenda et cæremoniarum, ritus, uerum in hijs etiam quæ credenda sunt.

As the Church of Hierusalem, Alexandria, and Antioche haue erred: so also the Church of Rome hath erred, not only in their liuing and maner of ceremonies, but also in matters of fayth.

As the Church of Jerusalem, Alexandria, and Antioch, have erred; so also the Church of Rome hath erred, not only in their living and manner of Ceremonies, but also in matters of Faith.

\xsection{Article XX. Of the Authority of the Church}

Habet Ecclesia Ritus statuendi ius, et in fidei controuersijs autoritatem, quamuis Ecclesiæ non licet quicquam instituere, quod verbo Dei scripto aduersetur, nec unum scripturæ locum sic exponere potest, ut alteri contradicat. Quare licet Ecclesia sit diuinorum librorum testis et conseruatrix, attamen vt aduersus eos nihil decernere, ita præter illos nihil credendum de necessitate salutis debet obtrudere.

The Church hath power to decree Rites or Ceremonies, and aucthoritie in controuersies of fayth: And yet it is not lawfull for the Church to ordayne any thyng that is contrarie to Gods worde written, neyther may it so expounde one place of scripture, that it be repugnaunt to another. Wherefore, although the Churche be a witnesse and a keper of holy writ: yet, as it ought not to decree any thing agaynst the same, so besides the same, ought it not to enforce any thing to be beleued for necessitie of saluation.

The Church hath power to decree Rites or Ceremonies, and authority in Controversies of Faith: and yet it is not lawful for the Church to ordain any thing that is contrary to God's Word written, neither may it so expound one place of Scripture, that it be repugnant to another. Wherefore, although the Church be a witness and a keeper of Holy Writ, yet, as it ought not to decree any thing against the same, so besides the same ought it not to enforce any thing to be believed for necessity of Salvation.

\xsection{Article XXI. Of the aucthoritie of generall Counselles}

\emph{Of the Authority of General Councils.}\footnote{The Twenty-first of the English Articles is omitted in the Amer. ed., because it is partly of a local and civil nature, and is provided for, as to the remaining parts of it, in other Articles.}

Generalia Concilia sine iussu et uoluntate principum congregari non possunt, et vbi conuenerint, quia ex hominibus constant, qui non omnes spiritu et uerbis Dei reguntur, et errare possunt, et interdum errarunt, etiam in hijs quæ ad normam pietatis pertinent: ideo quæ ab illis constituuntur,

Generall Counsels may not he gathered together without the commaundement and wyll of princes. And when they be gathered together (forasmuche as they be an assemblie of men, whereof all be not gouerned with the spirite and word of God) they may erre, and sometyme haue erred, euen in thinges parteynyng

ut ad salutem, necessaria, neque robur habent, neque autoritatem, nisi ostendi possint è sacris literis esse desumpta.

vnto God. Wherfore, thinges ordayned by them as necessary to saluation, haue neyther strength nor aucthoritie, vnlesse it may it declared that they be taken out of holy Scripture.

\xsection{Article XXII. Of Purgatory}

Doctrina Romanensium de Purgatorio, de Indulgentijs, de veneratione et adoratione tum Imaginum tum Reliquiarum, nec non de inuocatione Sanctorum, res est futilis, inaniter conficta, et nullis Scripturarum testimonijs innititur, imo verbo Dei contradicit.

The Romishe doctrine concernyng purgatorie, pardons, worshipping and adoration, as well of images, as of reliques, and also inuocation of Saintes, is a fonde thing, vainly inuented, and grounded vpon no warrantie of Scripture, but rather repugnaunt to the worde of God.

The Romish Doctrine concerning Purgatory, Pardons, Worshipping and Adoration, as well of Images as of Relics, and also Invocation of Saints, is a fond thing, vainly invented, and grounded upon no warranty of Scripture, but rather repugnant to the Word of God.

\xsection{Article XXIII. Of Ministering in the Congregation}

Non licet cuiquam sumere sibi munus publicè prædicandi, aut administrandi Sacramenta, in Ecclesia, nisi prius fuerit ad hæc obeunda legitimè uocatus et missus. Atque illos, legitimè uocatos et missos existimare debemus, qui per homines, quibus potestas uocandi Ministros atque mittendi in uineam Domini publicè concessa est in Ecclesia, cooptati fuerint et asciti in hoc opus.

It is not lawful for any man to take vpon hym the office of publique preachyng, or ministring the Sacramentes in the congregation, before he be lawfully called and sent to execute the same. And those we ought to iudge lawfully called and sent, whiche be chosen and called to this worke by men who haue publique aucthoritie geuen vnto them in the congregation, to call and sende ministers into the Lordes vineyarde.

It is not lawful for any man to take upon him the office of public preaching, or ministering the Sacraments in the Congregation, before he be lawfully called, and sent to execute the same. And those we ought to judge lawfully called and sent, which be chosen and called to this work by men who have public authority given unto them in the Congregation, to call and send Ministers into the Lord's vineyard.

\xsection{Article XXIV. Of Speaking in the Congregation in such a Tongue as the people understandeth}

Lingua populo non intellecta publicas in ecclesia preces peragere, aut Sacramenta administrare, verbo Dei et primitiuæ Ecclesiæ consuetudini planè repugnat.

It is a thing playnely repugnaunt to the worde of God, and the custome of the primitiue Churche, to haue publique prayer in the Churche, or to minister the Sacramentes in a tongue not vnderstanded of the people.

It is a thing plainly repugnant to the Word of God, and the custom of the Primitive Church, to have public Prayer in the Church, or to minister the Sacraments, in a tongue not understanded of the people.

\xsection{Article XXV. Of the Sacraments}

Sacramenta à Christo instituta, non tantum sunt notæ professionis Christianorum, sed certa quædam potius testimonia, et efficacia signa gratiæ, atque bonæ in nos uoluntatis Dei, per quæ inuisibiliter ipse in nobis operatur, nostrámque fidem in se, non solum excitat, uerumetiam confirmat.

Sacramentes ordayned of Christe, be not onely badges or tokens of Christian mens profession: but rather they be certaine sure witnesses and effectuall signes of grace and Gods good wyll towardes vs, by the which he doth worke inuisiblie in vs, and doth not only quicken, but also strengthen and confirme our fayth in hym.

Sacraments ordained of Christ be not only badges or tokens of Christian men's profession, but rather they be certain sure witnesses, and effectual signs of grace, and God's good will towards us, by the which he doth work invisibly in us, and doth not only quicken, but also strengthen and confirm our Faith in him.

Duo à Christo Domino nostro in Euangelio instituta sunt Sacramenta, scilicet Baptismus et Cœna Domini.

There are two Sacramentes ordayned of Christe our Lorde in the Gospell, that is to say, Baptisme, and the Supper of the Lorde.

There are two Sacraments ordained of Christ our Lord in the Gospel, that is to say, Baptism, and the Supper of the Lord.

Quinque illa uulgo nominata Sacramenta, scilicet, Confirmatio, Pœnitentia, Ordo, Matrimonium, et Extrema unctio, pro sacramentis euangelicis habenda non sunt, ut quæ partim à praua Apostolorum imitatione profluxerunt,

Those fyue, commonly called Sacramentes, that is to say, Confirmation, Penaunce, Orders, Matrimonie, and extreme Vnction, are not to be compted, for Sacramentes of the gospel, being such, as haue growen

Those five commonly called Sacraments, that is to say, Confirmation, Penance, Orders, Matrimony, and Extreme Unction, are not to be counted for Sacraments of the Gospel, being such as have grown

partim uitæ status sunt in scripturis quidem probati, sed sacramentorum eandem cum baptismo et cœna Domini rationem non habentes: quomodo nec Pœnitentia, ut quæ signum aliquod uisibile seu cæremoniam a Deo institutam non habeat.

partly of the corrupt folowing of the Apostles, partly are states of life alowed in the scriptures: but yet haue not lyke nature of Sacramentes with Baptisme and the Lordes Supper, for that they haue not any visible signe or ceremonie ordayned of God.

partly of the corrupt following of the Apostles, partly are states of life allowed in the Scriptures; but yet have not like nature of Sacraments with Baptism, and the Lord's Supper, for that they have not any visible sign or ceremony ordained of God.

\emph{Sacramenta, non in hoc instituta sunt à Christo, ut spectarentur, aut circumferentur, sed ut ritè illis uteremur: et in hijs duntaxat qui dignè percipiunt, salutarem habent effectum: qui uerò indigne percipiunt, damnationern} (\emph{ut inquit Paulus}) \emph{sibi ipsis acquirunt.}

The Sacramentes were not ordayned of Christ to be gased vpon, or to be caryed about; but that we shoulde duely use them. And in such only, as worthyly receaue the same, they haue a wholesome effect or operation: But they that receaue them vnworthyly, purchase to them selues damnation, as S. Paul sayth.

The Sacraments were not ordained of Christ to be gazed upon, or to be carried about, but that we should duly use them. And in such only as worthily receive the same, they have a wholesome effect or operation: but they that receive them unworthily, purchase to themselves damnation, as Saint Paul saith.

\xsection{Article XXVI. Of the Unworthiness of the Ministers, which hinders not the effect of the Sacraments}

Qvamuis in Ecclesia uisibili bonis mali semper sint admixti, atque interdum ministerio uerbi et sacramentorum administrationi præsint, tamen cùm non suo sed Christi nomine agant, eiúsque mandato et autoritate ministrent, illorum ministerio uti licet, cum in verbo Dei audiendo, tum in sacramentis percipiendis. Neque per illorum malitiam effectus

Although in the visible Churche the euyl be euer myngled with the good, and sometime the euyll haue cheefe aucthoritie in the ministration of the worde and Sacramentes: yet forasmuch as they do not the same in their own name but in Christes, and do minister by his commission and aucthoritie, we may vse their ministrie,

Although in the visible Church the evil be ever mingled with the good, and sometimes the evil have chief authority in the Ministration of the Word and Sacraments, yet forasmuch as they do not the same in their own name, but in Christ's, and do minister by his commission and authority, we may use their Ministry, both in

institutorum Christi tollitur, aut gratia donorum Dei minuitur, quoad eos qui fide et ritè sibi oblata percipiunt, quæ propter institutionem CHRISTI et promissionem efficacia sunt, licet per malos administrentur.

both in hearing the word of God, and in \emph{the} receauing \emph{of} the Sacramentes. Neither is y\textsuperscript{e} effect of Christes ordinaunce taken away by their wickednesse, nor the grace of Gods gyftes diminished from such as by fayth and ryghtly do receaue the Sacramentes ministered vnto them, which be effectuall, because of Christes institution and promise, although they be ministred by euyll men.

hearing the Word of God, and in receiving the Sacraments. Neither is the effect of Christ's ordinance taken away by their wickedness, nor the grace of God's gifts diminished from such as by faith, and rightly, do receive the Sacraments ministered unto them; which be effectual, because of Christ's institution and promise, although they be ministered by evil men.

Ad Ecclesiæ tamen disciplinam pertinent, ut in malos ministros inquiratur, accusentúrque ab hijs, qui eorum flagitia nouerint, atque tandem iusto conuicti iudicio, deponantur.

Neuerthelesse, it apperteyneth to the discipline of the Churche, that enquirie be made of euyl ministres, and that they be accused by those that haue knowledge of their offences: and finally, beyng founde gyltie by iust iudgement, be deposed.

Nevertheless, it appertaineth to the discipline of the Church, that inquiry be made of evil Ministers, and that they be accused by those that have knowledge of their offences; and finally, being found guilty, by just judgment be deposed.

\xsection{Article XXVII. Of Baptism}

Baptismus non est tantum professionis signum ac discriminis nota, qua Christiani à non Christianis discernantur, sed etiam est signum Regenerationis, per quod tanquam per instrumentum rectè baptismum suspitientes, ecclesiæ inseruntur, promissiones de Remissione peccatorum atque Adoptione nostra in filios Dei, per Spiritum sanctum uisibiliter

Baptisme is not onely a signe of profession, and marke of difference, whereby Christian men are discerned from other that be not christened: but is also a signe of regeneration or newe byrth, whereby as by an instrument, they that receaue baptisme rightly, are grafted into the Church: the promises of the forgeuenesse of sinne, and of

Baptism is not only a sign of profession, and mark of difference, whereby Christian men are discerned from others that be not christened, but \emph{it} is also a sign of Regeneration or New-Birth, whereby, as by an instrument, they that receive Baptism rightly are grafted into the Church; the promises of the forgiveness of sin, and

obsignantur, fides confirmatur, et ui diuinæ inuocationis, gratia augetur.

our adoption to be the sonnes of God, by the holy ghost, are visibly signed and sealed: fayth is confyrmed: and grace increased by vertue of prayer vnto God.

of our adoption to be the sons of God by the Holy Ghost, are visibly signed and sealed; Faith is confirmed, and Grace increased by virtue of prayer unto God.

Baptismus paruulorum omnino in ecclesia retinendus est, ut qui cum Christi institutione optimè congruat.

The baptisme of young children, is in any wyse to be retayned in the Churche, as most agreable with the institution of Christe.

The Baptism of young Children is in any wise to be retained in the Church, as most agreeable with the institution of Christ.

\xsection{Article XXVIII. Of the Lord's Supper}

Cœna Domini non est tantùm signum mutuæ beneuolentiæ Christianorum inter sese, uerum potiùs est sacramentum nostræ per mortem Christi redemptionis. Atque adeo ritè, dignè et cum fide sumentibus, panis quem frangimus, est communicatio corporis Christi: similiter poculum benedictionis, est communicatio sanguinis Christi.

The Supper of the Lord, is not only a signe of the loue that Christians ought to haue among them selues one to another: but rather it is a Sacrament of our redemption by Christes death. Insomuch that to suche as ryghtlie, worthyly, and with fayth receaue the same the bread whiche we breake is a parttakyng of the body of Christe, and likewyse the cuppe of blessing, is a parttakyng of the blood of Christe.

The Supper of the Lord is not only a sign of the love that Christians ought to have among themselves one to another; but rather it is a Sacrament of our Redemption by Christ's death: insomuch that to such as rightly, worthily, and with faith, receive the same, the Bread which we break is a partaking of the Body of Christ; and likewise the Cup of Blessing is a partaking of the Blood of Christ.

Panis et vini transubstantiatio in Eucharistia, ex sacris literis probari non potest, sed apertis scripturæ verbis aduersatur, sacramenti naturam euertit, et

Transubstantiation (or the chaunge of the substaunce of bread and wine) in the Supper of the Lorde, can not be proued by holye writ, but is repugnaunt to the playne wordes of scripture, ouerthroweth the nature of a Sacrament, and

Transubstantiation (or the change of the substance of Bread and Wine) in the Supper of the Lord, can not be proved by Holy Writ; but is repugnant to the plain words of Scripture, overthroweth the nature of a Sacrament, and

\emph{multarum superstitionum dedit occasionem}.\footnote{The following clause against the real presence and ubiquity of Christ's body was here added in the Parker Latin MS., but struck out in the Synod: '\emph{Christus in cœlum ascendens, corpori suo Immortalitatem dedit, Naturam non abstulit humane enim nature veritatem} (\emph{iuxta Scripturas}), \emph{perpetuo retinet, quam uno et definito Loco esse, et non in multa, vel omnia simul loca diffundi oportet. Quum igitur Christus in celum sublatus, ibi usque ad finem seculi permansurus, atque inde, non aliunde} (\emph{ut loquitur Augustinus}) \emph{venturus sit, ad iudicandum viuos et mortuos, non debet quisquam fidelium, et carnis eius, et sanguinis, realem, et corporalem} (\emph{ut loquuntur}) \emph{presentiam in Eucharistia vel credere, vel profiteri. Corpus tamen Christi datur,}' etc.}

hath geuen occasion to many superstitions.

hath given occasion to many superstitions.

Corpus Christi datur, accipitur, et manducatur in cœna, tantùm cœlesti et spirituali ratione. Medium autem quo Corpus Christi accipitur et manducatur in cœna, fides est.

The body of Christe is geuen, taken, and eaten in the Supper only after an heauenly and spirituall maner: And the meane whereby the body of Christe is receaued and eaten in the Supper, is fayth.

The Body of Christ is given, taken, and eaten, in the Supper, only after an heavenly and spiritual manner. And the mean whereby the Body of Christ is received and eaten in the Supper, is Faith.

Sacramentum Eucharistiæ ex institutione Christi non seruabatur, circumferebatur, eleuabatur, nec adorabatur.

The Sacrament of the Lordes Supper was not by Christes ordinaunce reserued, caryed about, lyfted vp, or worshipped.

The Sacrament of the Lord's Supper was not by Christ's ordinance reserved, carried about, lifted up, or worshiped.

[XXIX.\footnote{This Article, which agrees with the Zwinglian and Calvinistic theory against the Lutheran, is wanting in all the printed copies until 1571, and has here been supplied from the Parker MS. See Hardwick, p. 315, note 3, and p. 143.}

\xsection{Article XXIX. Of the Wicked, which eat not the Body of Christ in the use of the Lord's Supper}

[\emph{Impii, et fide viua destituti, licet carnaliter et visibiliter} (\emph{vt Augustinus loquitur}) \emph{corporis et sanguinis Christi sacramentum dentibus premant, nullo tamen modo Christi participes efficiuntur. Sed potius tantæ rei sacramentum seu symbolum}

The wicked, and suche as be voyde of a liuelye fayth, although they do carnally and visibly presse with their teeth (as Saint Augustine sayth) the Sacrament of the body and blood of Christ: yet in no wyse are the partakers of

The Wicked, and such as be void of a lively faith, although they do carnally and visibly press with their teeth (as Saint Augustine saith) the Sacrament of the Body and Blood of Christ; yet in no wise are they partakers of Christ:

adjudicium sibi mandueant et bibunt.]

Christe, but rather to their condemnation do eate and drinke the signe or Sacrament of so great a thing.

but rather, to their condemnation, do eat and drink the sign or Sacrament of so great a thing.

XXIX. [XXX.]

\xsection{Article XXX. Of both Kinds}

Calix Domini Laicis non est denegandus: utraque enim pars dominici sacramenti ex Christi institutione et præcepto, omnibus Christianis ex æquo administrari debet.

The cuppe of the Lorde is not to be denyed to the laye people. For both the partes of the Lordes Sacrament, by Christes ordinaunce and commaundement, ought to be ministred to all Christian men alike.

The Cup of the Lord is not to be denied to the Lay-people: for both the parts of the Lord's Sacrament, by Christ's ordinance and commandment, ought to be ministered to all Christian men alike.

XXX. [XXXI.]

\xsection{Article XXXI. Of the one Oblation of Christ finished upon the Cross}

\emph{Oblatio Christi semel facta, perfecta est redemptio, propitiatio, et satisfactio pro omnibus peccatis totius mundi, tam originalibus quam actualibus. Neque præter illam unicam est ulla alia pro peccatis expiatio. Vnde missarum sacrificia, quibus uulgo dicebatur, Sacerdotem offerre Christum in remissionem pœna aut culpæ pro uiuis et defunctis, blasphema figmenta sunt, et pernitiosæ imposturæ.}

The offering of Christ once made, is \emph{the} perfect redemption, propiciation, and satisfaction for all the sinnes of the whole worlde, both originall and actuall, and there is none other satisfaction for sinne, but that alone. Wherefore the sacrifices of Masses, in the which it was commonly said that the Priestes did offer Christe for the quicke and the dead, to haue remission of payne or gylt, were blasphemous fables, and daungerous deceits.

The Offering of Christ once made is \emph{that} perfect redemption, propitiation, and satisfaction, for all the sins of the whole world, both original and actual; and there is none other satisfaction for sin, but that alone. Wherefore the sacrifices of Masses, in the which it was commonly said that the Priest did offer Christ for the quick and the dead, to have remission of pain or guilt, were blasphemous fables, and dangerous deceits.

XXXI. [XXXII.]

\xsection{Article XXXII. Of the Marriage of Priests}

Episcopis, Presbyteris, et Diaconis, nullo mandato diuino

Byshops, Priestes, and Deacons, are not commaunded

Bishops, Priests, and Deacons, are not commanded

præceptum est, ut aut cœlibatum uoueant, aut à matrimonio abstineant. Licet igitur etiam illis, vt cœteris omnibus Christianis, vbi hoc ad pietatem magis facere iudicauerint, pro suo arbitratu matrimonium contrahere.

by Gods lawe eyther to vowe the estate of single lyfe, or to abstayne from mariage. Therefore it is lawfull \emph{also} for them, as for all other Christian men, to mary at ther owne discretion, as they shall iudge the same to serue better to godlynesse.

by God's Law, either to vow the estate of single life, or to abstain from marriage: therefore it is lawful for them, as for all other Christian men, to marry at their own discretion, as they shall judge the same to serve better to godliness.

XXXII. [XXXIII.]

\xsection{Article XXXIII. Of excommunicate Persons, how they are to be avoided}

Qvi per publicam Ecclesiæ denuntiationem ritè ab unitate ecclesiæ præcisus est et excommunicatus, is ab uniuersa fidelium multitudine, donec per pœnitentiam publicè reconciliatus fuerit, arbitrio Iudicis competentis, habendus est tanquam Ethnicus et Publicanus.

That person whiche by open denuntiation of the Churche, is rightly cut of from the vnitie of the Churche, and excommunicated, ought to be taken of the whole multitude of the faythfull as an Heathen and Publicane, vntill he be openly reconciled by penaunce, and receaued into the Churche by a iudge that hath aucthoritie \emph{thereto}.

That person which by open denunciation of the Church is rightly cut off from the unity of the Church, and excommunicated, ought to be taken of the whole multitude of the faithful, as a Heathen and Publican, until he be openly reconciled by penance, and received into the Church by a judge that hath authority \emph{thereunto.}

XXXIII. [XXXIV.]

\xsection{Article XXXIV. Of the Traditions of the Church}

Traditiones atque cæremonias easdem, non omnino necessarium est esse ubique aut prorsus consimiles. Nam et uariæ semper fuerunt, et mutari possunt, pro regionum, temporum, et morum diuersitate, modo nihil contra uerbum Dei instituatur.

It is not necessarie that traditions and ceremonies be in al places one, or vtterly like, for at all times they haue ben diuerse, and may be chaunged accordyng to the diuersitie of countreys, times, and mens maners, so that nothing be ordeyned against Gods worde.

It is not necessary that Traditions and Ceremonies be in all places one, or utterly like; for at all times they have been divers, and may be changed according to the diversity of countries, times, and men's manners, so that nothing be ordained against God's Word.

Traditiones et cæremonias ecclesiasticas quæ cum uerbo Dei non pugnant, et sunt autoritate publica institutæ atque probatæ, quisquis priuato consilio uolens et data opera publicè uiolauerit, is, ut qui peccat in publicum ordinem ecclesiæ, quique lædit autoritatem Magistratus, et qui infirmorum fratrum conscientias uulnerat, publicè, ut cœteri timeant, arguendus est.

Whosoeuer through his priuate iudgement, wyllyngly and purposely doth openly breake the traditions and ceremonies of the Church, which be not repugnaunt to the worde of God, and be ordayned and approued by common aucthoritie, ought to be rebuked openly (that other may feare to do the lyke), as he that offendeth agaynst the Common order of the Churche and hurteth the aucthoritie of the Magistrate, and woundeth the consciences of the weake brethren.

Whosoever, through his private judgment, willingly and purposely, doth openly break the Traditions and Ceremonies of the Church, which be not repugnant to the Word of God, and be ordained and approved by common authority, ought to be rebuked openly (that others may fear to do the like), as he that offendeth against the common order of the Church, and hurteth the authority of the Magistrate, and woundeth the consciences of the weak brethren.

Quæ libet ecclesia particularis, sine nationalis autoritatem habet instituendi, mutandi, aut abrogandi cæremonias aut ritus Ecclesiasticos, humana tantum autoritate institutos, modò omnia ædificationem fiant.

Euery particuler or nationall Churche, hath aucthoritie to ordaine, chaunge, and abolishe ceremonies or rites of the Churche ordeyned onlye by mans aucthoritie, so that all thinges be done to edifiyng.

Every particular or national Church hath authority to ordain, change, and abolish, Ceremonies or Rites of the Church ordained only by man's authority, so that all things be done to edifying.

XXXIV. [XXXV.]

\xsection{Article XXXV. Of the Homilies}

Tomus secundus Homiliarum, quarum singulos titulos huic Articulo subiunximus, continet piam et salutarem doctrinam, et hijs temporibus necessariam, non minus quàm prior Tomus Homiliarum quæ æditæ sunt tempore Edwardi sexti. Itaque eas in ecclesijs per ministros diligenter et clarè, ut à populo

The seconde booke of Homilies, the seuerall titles whereof we haue ioyned vnder this article, doth conteyne a godly and wholesome doctrine, and necessarie for these tymes, as doth the former booke of Homilies, which were set foorth in the time of Edwarde the sixt: and

The Second Book of Homilies, the several titles whereof we have joined under this Article, doth contain a godly and wholesome Doctrine, and necessary for these times, as doth the former Book of Homilies, which were set forth in the time of Edward the Sixth; and therefore

intelligi possint, recitandas, esse iudicamus.

therefore we iudge them to be read in Churches by the Ministers diligently, and distinctly, that they may be vnderstanded by the people.

we judge them to be read in Churches by the Ministers, diligently and distinctly, that they may be understanded of the people.

[XXXIV.]

Catalogus Homiliarum.

Of the names of the Homilies.

Of the names of the Homilies.

De recto ecclesiæ usu.

Aduersus Idololatriæ pericula.

De reparandis ac purgandis ecclesijs.

De bonis operibus.

De ieiunio.

In gulæ atque ebrietatis uitia.

In nimis sumptuosos uestium apparatus.

De oratione siue precatione.

De loco et tempore orationi destinatis.

De publicis precibus ac Sacramentis, idiomate uulgari omnibusque noto, habendis.

De sacrosancta uerbi divini autoritate.

De eleemosina.

De Christi natiuitate.

De dominica passione.

De resurrectione Domini.

De digna corporis et sanguinis dominici in cœna Domini participatione.

De donis spiritus sancti.

In diebus, qui uulgo Rogationum dicti sunt, concio.

De matrimonij statu.

De otio seu socordia.

De pœnitentia.

1 Of the right vse of the Churche.

2 Agaynst perill of Idolatrie.

3 Of repayring and keping cleane of Churches.

4 Of good workes, first of fastyng.

5 Agaynst gluttony and drunkennesse.

6 Agaynst excesse of apparell.

7 Of prayer.

8 Of the place and time of prayer.

9 That common prayer and Sacramentes ought to be ministred in a knowen tongue.

10 Of the reuerente estimation of Gods worde.

11 Of almes doing.

12 Of the Natiuitie of Christe.

13 Of the passion of Christe.

14 Of the resurrection of Christe.

15 Of the worthie receauing of the Sacrament of the body and blood of Christe.

16 Of the gyftes of the holy ghost.

17 For the Rogation dayes.

1. Of the right Use of the Church.

2. Against Peril of Idolatry.

3. Of repairing and keeping clean of Churches.

4. Of good Works: first of Fasting.

5. Against Gluttony and Drunkenness.

6. Against Excess of Apparel.

7. Of Prayer.

8. Of the Place and Time of Prayer.

9. That Common Prayers and Sacraments ought to be ministered in a known tongue.

10. Of the reverend Estimation of God's Word.

11. Of Alms-doing.

12. Of the Nativity of Christ.

13. Of the Passion of Christ.

14. Of the Resurrection of Christ.

15. Of the worthy receiving of the Sacrament of the Body and Blood of Christ.

16. Of the Gifts of the Holy Ghost.

17. For the Rogation-days.

18 Of the state of Matrimonie.

19 Of repentaunce.

20 Agaynst Idlenesse.

21 Agaynst rebellion.

18. Of the State of Matrimony.

19. Of Repentance.

20. Against Idleness.

21. Against Rebellion.

[\emph{This Article is received in this Church, so far as it declares the Books of Homilies to be an explication of Christian doctrine, and instructive in piety and morals. But all references to the constitution and laws of England are considered as inapplicable to the circumstances of this Church; which also suspends the order for the reading of said Homilies in churches, until a revision of them may be conveniently made, for the clearing of them, as well from obsolete words and phrases, as from, the local references.}]

XXXV. [XXXVI.]

\xsection{Article XXXVI. Of Consecration of Bishops and Ministers}

Libellus de Consecratione Archiepiscoporum \& Episcoporum, \& de ordinatione Presbyterorum \& Diaconorum æditus nuper temporibus Edwardi sexti, \& autoritate Parlamenti illis ipsis temporibus confirmatus, omnia ad eiusmodi consecrationem \& ordinationem necessaria continet, \& nihil habet quod ex se sit aut superstitiosum aut impium.

The booke of Consecration of \emph{Archbishops, and} Byshops, and orderyng of Priestes and Deacons, \emph{lately set foorth in the time of Edwarde the sixt, and confirmed at the same tyme by aucthoritie of Parliament}, doth conteyne all thinges necessarie to suche consecration and orderyng: neyther hath it any thing, that of it selfe is superstitious

The Book of Consecration of Bishops, and Ordering of Priests and Deacons, as set forth by the General Convention of this Church in 1792, doth contain all things necessary to such Consecration and Ordering; neither hath it any thing that, of itself, is superstitious and ungodly. And, therefore, whosoever are consecrated or ordered

Itaque quicunque iuxta ritus illius libri consecrati aut ordinati sunt ab Anno secundo prædicti Regis Edwardi, usque ad hoc tempus, aut in posterum iuxta eosdem ritus consecrabuntur aut ordinabuntur ritè, ordine, atque legitimè, statuimus esse \& fore consecratos \& ordinatos.

\emph{or} vngodly. And therefore, whosoeuer are consecrate or ordered accordyng \emph{to the rites of that booke, sence the seconds yere of the aforenamed king Edwarde, vnto this time or hereafter shal be consecrated or ordered accordyng to the same rites}, we decree all such to be ryghtly, orderly, and lawfully consecrated and ordered.

according to \emph{said Form}, we decree all such to be rightly, orderly, and lawfully consecrated and ordered.

XXXVI. [XXXVII.]

\xsection{Article XXXVII. Of the Power of the Civil Magistrates}

Regia Maiestas in hoc Angliæ Regno ac cœteris eius Dominijs, iure summam habet potestatem, ad quam omnium statuum huius Regni, siue illi ecclesiastici sunt siue non, in omnibus causis suprema gubernatio pertinet, \& nulli externæ iurisdictioni est subiecta, nec esse debet.

The Queenes Maiestie hath the cheefe power in this Realme of Englande, and other her dominions, vnto whom the cheefe gouernment of all estates of this Realme, whether they be Ecclesiasticall or Ciuile, in all causes doth apparteine, and is not, nor ought to be subiect to any forraigne iurisdiction.

The Power of the Civil Magistrate extendeth to all men, as well Clergy as Laity, in all things temporal; but hath no authority in things purely spiritual. And, we hold it to be the duty of all men who are professors of the Gospel, to pay respectful obedience to the Civil Authority, regularly and legitimately constituted.

Cum Regiæ Maiestati summam gubernationem tribuimus,. quibus titulis intelligimus animos quorundam calumniatorum offendi: non damus Regibus nostris aut uerbi Dei aut sacramentorum administrationem, quod etiam Iniunctiones ab Elizabetha Regina nostra nuper æditæ, apertissimè testantur: sed eam tantùm prærogatium, quam in sacris scripturis à Deo ipso omnibus pijs

Where we attribute to the Queenes Maiestie the cheefe gouernment, by whiche titles we vnderstande the mindes of some slanderous folkes to be offended: we geue not to our princes the ministring either of God's word, or of Sacraments, the which thing the injunctions also lately set forth by Elizabeth our Queene, doth most plainlie testifie: But that only prerogatiue whiche we see to

Principibus, uidemus semper fuisse attributam, hoc est, ut omnes status atque ordines fidei suæ a Deo commissos, siue illi ecclesiastici sint, siue ciuiles, in officio contineant, \& contumaces ac delinquentes, gladio ciuili coerceant.

haue ben geuen alwayes to all godly Princes in holy Scriptures by God him selfe, that is, that they should rule all estates and degrees committed to their charge by God, whether they be Ecclesiasticall or Temporall, and restraine with the ciuill sworde the stubberne and euyll doers.

Romanus Pontifex nullam habet iurisdictionem in hoc regno Angliæ.

The bishop of Rome hath no iurisdiction in this Realme of Englande.

Leges Ciuiles possunt Christianas propter capitalia et grauia crimina morte punire.

The lawes of the Realme may punishe Christian men with death, for heynous and greeuous offences.

Christianis licet et ex mandato Magistratus arma portare, et iusta bella administrare.

It is lawfull for Christian men, at the commaundement of the Magistrate, to weare weapons, and serue in the warres.

XXXVII. [XXXVIII.]

\xsection{Article XXXVIII. Of Christian Men's Goods, which are not common}

Facultates \& bona Christianorum non sunt communia quoad ius \& possessionem, vt quidam Anabaptistæ falso iactant. Debet tamen quisque de hijs quæ possidet, pro facultatum ratione, pauperibus eleemosynas benigne distribure.

The ryches and goodes of Christians are not common, as touching the ryght, title, and possession of the same, as certayne Anabaptistes do falsely boast. Notwithstandyng euery man ought of suche thinges as he possesseth, liberally to geue almes to the poore, accordyng to his habilitie.

The Riches and Goods of Christians are not common, as touching the right, title, and possession of the same; as certain Anabaptists do falsely boast. Notwithstanding, every man ought, of such things as he possesseth, liberally to give alms to the poor, according to his ability.

XXXVIII. [XXXIX.]

\xsection{Article XXXIX. Of a Christian Man's Oath}

Qvemadmodum iuramentum

As we confesse that vayne

As we confess that vain

uanum \& temerarium à Domino nostro Iesu Christo, \& Apostolo eius Iacobo, Christianis hominibus interdictum esse fatemur: ita Christianam religionem minimè prohibere censemus, quin iubente Magistratu, in causa fidei \& charitatis, iurare liceat, modò id fiat iuxta, Prophetæ doctrinam, in iustitia, in iudicio, \& ueritate.\footnote{In the Forty-two Articles of Edward VI. there are four additional Articles—on the Resurrection of the Dead, the State of the Souls of the Departed, Millenarians, and the Eternal Damnation of the Wicked.}

and rashe swearing is forbidden Christian men by our Lord Jesus Christe, and James his Apostle: So we iudge that Christian religion doth not prohibite, but that a man may sweare when the Magistrate requireth, in a cause of faith and charitie, so it be done accordyng \emph{to the prophetes} teaching, in iustice, iudgement, and trueth.\footnote{In the Forty-two Articles of Edward VI. there are four additional Articles—on the Resurrection of the Dead, the State of the Souls of the Departed, Millenarians, and the Eternal Damnation of the Wicked.}

and rash Swearing is forbidden Christian men by our Lord Jesus Christ, and James his Apostle, so we judge, that Christian Religion doth not prohibit, but that a man may swear when the Magistrate requireth, in a cause of faith and charity, so it be done according \emph{to the Prophets'} teaching, in justice, judgment, and truth.

[The remainder of the English editions is omitted in the American Revision.]

The Ratification.

Hos Articulos fidei Christianæ, continentes in uniuersum nouemde cimpaginas in autographo, quod asseruatur apud Reuerendissimum in Christo patrem, Dominum Matthæum Centuariensem Archiepiscopum, totius Angliæ Primatem \& Metropolitanum, Archiepiscopi \& Episcopi utriusque Prouinciæ regni Angliæ, in sacra prouinciali Synodo legitimè congregati, unanimi assensu recipiunt \& profitentur, \& ut ueros atque Orthodoxos, manuum suarum subscriptionibus approbant, uicesimo nono die mensis Ianuarij: Anno Domini,

This Booke of Articles before rehearsed, is agayne approued, and allowed to be holden and executed within the Realme, by the ascent and consent of our Soueraigne Ladye Elizabeth, by the grace of GOD, of Englande, Fraunce, and Irelande Queene, defender of the fayth, \&c. Which Articles were deliberately read, and confirmed agayne by the subscription of the handes of the Archbyshop and Byshoppes of the vpper house, and by the subscription of the whole Cleargie in the neather house in their Conuocation,

secundum computationem ecclesiæ Anglicanæ, millesimo quingentesimo sexagesimo secundo: uniuersusque Clerus Inferioris domus, eosdem etiam unanimiter \& recepit \& professus est, ut ex manuum suarum subscriptionibus patet, quas obtulit \& deposuit apud eundem Reuerendissimum, quinto die Februarij, Anno prædicto.

Quibus omnibus articulis, Serenesima princeps Elizabeth, Dei gratia Angliæ, Franciæ, \& Hiberniæ Regina, fidei Defensor, \&c. per seipsam diligenter prius lectis \& examinatis, Regium suum assensum præbuit.

in the yere of our Lorde GOD. 1571.

[A Table of the Articles.\}\footnote{This heading is inserted in the later English editions after the Ratification.}

1 Of fayth in the Trinitie.

2 Of Christe the sonne of GOD.

3 Of his goyng downe into hell.

4 Of his Resurrection.

5 Of the holy ghost.

6 Of the sufficiencie of the Scripture.

7 Of the olde Testament.

8 Of the three Credes.

9 Of originall sinne.

10 Of free wyll.

11 Of Iustification.

12 Of good workes.

13 Of workes before iustification.

14 Of workes of supererogation.

15 Of Christe alone without sinne.

16 Of sinne after Baptisme.

17 Of predestination and election.

18 Of obtayning saluation by Christe.

19 Of the Churche.

20 Of the aucthoritie of the Churche.

21 Of the aucthoritie of generall Counsels.

22 Of Purgatorie.

23 Of ministring in the congregation.

24 Of speaking in the congregation.

25 Of the Sacramentes.

26 Of the vnworthynesse of the Ministers.

27 Of Baptisme.

28 Of the Lordes supper.

29 Of the wicked whiche eate not the body of Christe.

30 Of both kyndes.

31 Of Christes one oblation.

32 Of the mariage of Priestes.

33 Of excommunicate persons.

34 Of traditions of the Churche.

35 Of Homilies.

36 Of consecration of Ministers.

37 Of ciuill Magistrates.

38 Of Christian mens goods.

39 Of a Christian mans othe.

40 Of the ratification.

Excusum Londini apud \textsc{REGINALDVM} Wolfium, Regiæ Maiest. in Latinis typographum. \textsc{anno domini}. 1563.

¶ Imprinted at London in Powles Churchyard, by Richarde Iugge and Iohn Cawood, Printers to the Queenes Maiestie, in Anno Domini 1571.

* Cum priuilegio Regise maiestatis.
